\appendix

\section{Full Proofs}
\label{app:proofs}

\subsection{Proof of Theorem~\ref{thm:main}}
\label{app:main_proof}

We provide a complete proof of the regret bound. Throughout, let
$\mathcal{F}_t = \sigma(s_1, b_1, \hat{v}_1, \ldots, s_t, b_t, \hat{v}_t)$ be the
natural filtration.

\begin{lemma}[Proxy estimation error]
\label{lem:proxy}
Under Assumptions~\ref{ass:lipschitz} and \ref{ass:noise}, for any step $t$ with
$b_t = 0$, the proxy satisfies
\[
  \mathbb{E}\!\left[|V(s_t) - \hat{v}_t| \;\middle|\; \mathcal{F}_{t-1}\right]
  \;\leq\; \tau_t + \sigma_{\max}.
\]
\end{lemma}

\begin{proof}
When $b_t = 0$, the algorithm sets $\hat{v}_t = \hat{V}(s_{\text{anc}})$ where
$s_{\text{anc}}$ is the stored ancestor. We have
\begin{align*}
  |V(s_t) - \hat{v}_t|
  &\leq |V(s_t) - V(s_{\text{anc}})| + |V(s_{\text{anc}}) - \hat{V}(s_{\text{anc}})|\\
  &\leq L \cdot |d(s_t) - d(s_{\text{anc}})| + |\varepsilon_{s_{\text{anc}}}|.
\end{align*}
The condition $b_t = 0$ means $L \cdot |d(s_t) - d(s_{\text{anc}})| \leq \tau_t$.
Taking expectations and using $\mathbb{E}[|\varepsilon|] \leq \sigma_{\max}$ gives
the result.
\end{proof}

\begin{lemma}[Estimation error when verifier is called]
\label{lem:direct}
Under Assumption~\ref{ass:noise}, for any step $t$ with $b_t = 1$,
\[
  \mathbb{E}\!\left[|V(s_t) - \hat{v}_t| \;\middle|\; \mathcal{F}_{t-1}\right]
  \;\leq\; \sigma_{\max}.
\]
\end{lemma}

\begin{proof}
Immediate from $\hat{v}_t = \hat{V}(s_t) = V(s_t) + \varepsilon_{s_t}$ and the
sub-Gaussian bound on $\varepsilon_{s_t}$.
\end{proof}

\begin{lemma}[UCT selection error, depth-$D$ tree]
\label{lem:uct}
For a tree of branching factor $K$ and depth $D$, let
\[
  \Phi(K,D) \;=\; \sum_{d=0}^{D-1} K^{d/2} \;=\; \frac{K^{D/2}-1}{\sqrt{K}-1}
  \quad (K > 1).
\]
The cumulative selection error satisfies
\[
  \sum_{t=1}^{T} \mathbb{E}\!\left[\hat{v}_{s^*_t} - V(\hat{s}^*_t)\right]
  \;\leq\; 8\sigma_{\max}\sqrt{2} \cdot \Phi(K,D) \cdot \sqrt{T\log K}.
\]
\end{lemma}

\begin{proof}
\textbf{Single-node UCT bound.}
For any node $v$ visited $n_v \geq 1$ times, the same UCT analysis as in
the standard $K$-armed bandit setting \citep{auer2002finite} applies to its $K$
children: the selection overestimation at $v$ over its $n_v$ visits is bounded by
$8\sigma_{\max}\sqrt{2 n_v \log K}$.

We verify the argument.  Let $c_a(s) = \sigma_{\max}\sqrt{2\ln(n_v)/n_a(s)}$ be the UCT
confidence width for child $a$ after $n_a(s)$ visits within $v$'s sub-problem.
By Hoeffding's inequality for sub-Gaussian noise:
\[
  \Pr\!\left[|\hat{v}_a^{(s)} - V(a)| > c_a(s)\right] \leq \frac{2}{n_v^2},
\]
and a union bound over $K$ children and $n_v$ rounds gives high-probability validity
of all confidence intervals.  On this event, overestimation at each round contributes
at most the confidence width $c_{a}(s)$.  Summing over $n_v$ rounds by Cauchy--Schwarz
and absorbing logarithmic factors:
\[
  \mathcal{E}_v \;\leq\; 8\sigma_{\max}\sqrt{2 n_v \log K}.
\]

\textbf{Depth-$D$ extension via backward induction.}
Label depth levels $0$ (root) to $D-1$ (parents of leaves).  At each depth $d$, let
$\mathcal{V}_d$ be the set of internal nodes; each node $v \in \mathcal{V}_d$ is
visited $n_v$ times in total.  The total visits at depth $d$ sum to $T$:
$\sum_{v \in \mathcal{V}_d} n_v = T$.

Applying the single-node bound at every node at depth $d$ and using
Cauchy--Schwarz ($\sum_v \sqrt{n_v} \leq \sqrt{|\mathcal{V}_d|} \cdot \sqrt{T}$,
since $|\mathcal{V}_d| = K^d$):
\[
  \sum_{v \in \mathcal{V}_d} \mathcal{E}_v
  \;\leq\; \sum_{v \in \mathcal{V}_d} 8\sigma_{\max}\sqrt{2 n_v \log K}
  \;\leq\; 8\sigma_{\max}\sqrt{2\log K} \cdot K^{d/2} \cdot \sqrt{T}.
\]

Summing over all $D$ levels and recognising the geometric series:
\[
  \sum_{t=1}^T \mathbb{E}\!\left[\hat{v}_{s^*_t} - V(\hat{s}^*_t)\right]
  \;\leq\; \sum_{d=0}^{D-1} 8\sigma_{\max}\sqrt{2\log K} \cdot K^{d/2} \cdot \sqrt{T}
  \;=\; 8\sigma_{\max}\sqrt{2T\log K} \cdot \Phi(K,D). \qquad\square
\]
\end{proof}

\begin{remark}[Depth--constant tradeoff]
\label{rem:tree_depth}
For $K{=}2$, $D{=}12$: $\Phi(2,12) = \sum_{d=0}^{11}2^{d/2} \leq 153$, so the proof
gives $C_1 \leq 8 \cdot 153 \cdot \sigma_{\max}\sqrt{2} \approx 1224\,\sigma_{\max}\sqrt{2}$.
This constant is large but fixed (independent of $T$), confirming the
$O(\sqrt{T\log K})$ rate.  It is a Cauchy--Schwarz upper bound and typically
far exceeds the empirical value: the observed slope $0.635$ at $T{=}400$
(Figure~\ref{fig:regret}) and the actual regret values ($\approx 31$ at $T{=}400$)
are consistent with an effective constant in the range $1$--$5\,\sigma_{\max}\sqrt{2}$
for the paper's parameters.  The gap is expected: the Cauchy--Schwarz step is tight
only when all $K^d$ nodes at depth $d$ are visited equally, which does not occur in
practice under guided UCT exploration.
\end{remark}

\begin{proof}[Proof of Theorem~\ref{thm:main}]
Decompose regret at each step $t$ as
\begin{equation}
  V(s^*_t) - V(\hat{s}^*_t)
  = \underbrace{V(s^*_t) - \hat{v}_{s^*_t}}_{\text{(I) estimation error}}
  + \underbrace{\hat{v}_{s^*_t} - V(\hat{s}^*_t)}_{\text{(II) selection error}}.
  \label{eq:decomp}
\end{equation}

\paragraph{Bounding term (II): selection error.}
Summing over $t$ and applying Lemma~\ref{lem:uct}:
\[
  \sum_{t=1}^T \mathbb{E}\!\left[\hat{v}_{s^*_t} - V(\hat{s}^*_t)\right]
  \;\leq\; C_1 \sqrt{T\log K},
  \qquad C_1 = 8\sigma_{\max}\sqrt{2}\,\Phi(K,D).
\]

\paragraph{Bounding term (I): estimation error.}
Term~(I) is the \emph{signed} quantity $V(s^*_t) - \hat{v}_{s^*_t}$, not its
absolute value.  Under the zero-mean noise assumption (Assumption~\ref{ass:noise}),
the two cases are:

\emph{Case $b_t = 1$ (verifier called):}
$\hat{v}_{s^*_t} = \hat{V}(s^*_t) = V(s^*_t) + \varepsilon_{s^*_t}$, so
$\mathbb{E}\!\left[V(s^*_t) - \hat{v}_{s^*_t}\right] = \mathbb{E}[-\varepsilon_{s^*_t}] = 0$.

\emph{Case $b_t = 0$ (proxy used):}
$\hat{v}_{s^*_t} = \hat{V}(s_{\text{anc}}) = V(s_{\text{anc}}) + \varepsilon_{s_{\text{anc}}}$,
so
\[
  \mathbb{E}\!\left[V(s^*_t) - \hat{v}_{s^*_t}\right]
  = V(s^*_t) - V(s_{\text{anc}}) + \mathbb{E}[-\varepsilon_{s_{\text{anc}}}]
  = V(s^*_t) - V(s_{\text{anc}})
  \;\leq\; L\cdot|d(s^*_t) - d(s_{\text{anc}})|
  \;\leq\; \tau_t,
\]
where the last inequality uses the condition $b_t = 0$
(i.e.\ $L\cdot|d(s_t)-d(s_{\text{anc}})| \leq \tau_t$) and
Assumption~\ref{ass:lipschitz}.

Summing over all steps:
\begin{align}
  \sum_{t=1}^T \mathbb{E}\!\left[V(s^*_t) - \hat{v}_{s^*_t}\right]
  &= \sum_{t:\,b_t=1} 0
   + \sum_{t:\,b_t=0} \mathbb{E}\!\left[V(s^*_t) - \hat{v}_{s^*_t}\right]
   \notag\\
  &\leq \sum_{t:\,b_t=0} \tau_t
  \;\leq\; \gamma\sum_{t=1}^T t^{-1/2}
  \;\leq\; 2\gamma\sqrt{T}.
  \label{eq:est_err}
\end{align}
The zero-mean property eliminates the $\sigma_{\max}$ contribution from
direct-call steps; for proxy steps, the ancestor's noise cancels in expectation.

\paragraph{Combining.}
Adding term~(II) (Lemma~\ref{lem:uct}) and term~(I) (equation~\eqref{eq:est_err}),
using the shorthand $C_1 = C_1(K,D) = 8\sigma_{\max}\sqrt{2}\,\Phi(K,D)$:
\[
  \mathbb{E}[R(T)]
  \;\leq\; \underbrace{C_1\sqrt{T\log K}}_{\text{(II) selection}}
         + \underbrace{2\gamma\sqrt{T}}_{\text{(I) estimation}}.
\]

\paragraph{Optimising $\gamma$.}
Setting
\[
  \gamma^* = \tfrac{C_1}{2}\sqrt{\log K}
           = 4\sigma_{\max}\sqrt{2\log K}\cdot\Phi(K,D)
\]
gives $2\gamma^*\sqrt{T} = C_1\sqrt{T\log K}$, so both terms are equal and
\[
  \mathbb{E}[R(T)]
  \;\leq\; 2C_1\sqrt{T\log K}
  \;=\; C(K,D)\sqrt{T\log K},
  \quad C(K,D) = 2C_1 = 16\sigma_{\max}\sqrt{2}\,\Phi(K,D). \qquad\square
\]
\end{proof}

\subsection{Proof of Theorem~\ref{thm:calls}}
\label{app:calls}

We prove the bound in two parts corresponding to the exploration and exploitation
regimes identified in Remark~\ref{rem:exploration}.

\begin{lemma}[Single-verification property]
\label{lem:single_verify}
Algorithm~\ref{alg:dva} invokes the verifier on any node $s \in \mathcal{T}$
at most once.
\end{lemma}

\begin{proof}
Line~11 of Algorithm~\ref{alg:dva} stores the observation $(s', \hat{v}(s'))$ in
the score memory $\mathcal{M}$ immediately after calling the verifier.  At every
subsequent visit to $s'$, line~8 finds $s' \in \mathcal{M}$, so $s_{\text{anc}} = s'$
and $\delta = L\cdot 0 = 0 < \tau_t$ for any $\tau_t > 0$; the else branch is taken
and the verifier is not called.  (In the implementation this is enforced by the
\texttt{verifier\_called} flag on each node.)
\end{proof}

\begin{proof}[Proof of Theorem~\ref{thm:calls}]

\noindent\textbf{Part (a): Single-verification ceiling.}
By Lemma~\ref{lem:single_verify}, $\mathcal{B}(\pi) \leq |\mathcal{T}| \leq N_\mathcal{T}$,
where $N_\mathcal{T} = K(K^D-1)/(K-1) = K + K^2 + \cdots + K^D$ is the number of
non-root nodes (verifier calls can occur at depths $1$ through $D$ only, since the
root has no ancestor from which to compute a depth gap).  Since $K$ and $D$
are fixed problem parameters, $N_\mathcal{T}$ is a constant independent of $T$, and
the call fraction $\mathcal{B}/T \leq N_\mathcal{T}/T \to 0$ as $T \to \infty$.

\medskip
\noindent\textbf{Part (b): Threshold-triggered refinement.}
A \emph{threshold-triggered} call requires $s_\mathrm{anc} \neq \emptyset$
(re-visit) and $\delta_t = L \cdot |d(s_t) - d(s_\mathrm{anc})| > \tau_t$.
Since depths lie in $\{0,\ldots,D\}$, we have $\delta_t \leq LD$ always.
When $\tau_t \geq LD$---equivalently, $t \leq (\gamma/(LD))^2 =: T_\mathrm{free}$---the
threshold dominates every possible depth gap, so no threshold-triggered call
occurs.  Verifier calls during $t \leq T_\mathrm{free}$ are exclusively
first-visit calls ($s_\mathrm{anc} = \emptyset$), which are already bounded
in total by $N_\mathcal{T}$ via Part~(a).
The number of threshold-triggered calls is at most $\max(0, T - T_\mathrm{free})$;
together with Part~(a) the total satisfies $\mathcal{B}(\pi) \leq N_\mathcal{T}$.

\medskip
\noindent\textbf{Part (c): Savings fraction.}
Immediate from part (a): $\mathcal{B}(\pi)/T \leq N_\mathcal{T}/T \to 0$.

\medskip
\noindent\textbf{Empirical correspondence.}
For the paper's parameters ($K{=}2$, $D{=}12$, $N_\mathcal{T}{=}8{,}190$,
$\gamma{=}1.0$, $L{=}0.35$): $T_\mathrm{free} = (1.0/(0.35 \times 12))^2 \approx 0.06$,
so the call-free phase is negligible.  In the experimental range $T \leq 3200$
(exploration regime, $T \ll N_\mathcal{T}$), calls grow approximately as
$1.9\sqrt{T}\log T$ empirically.  The formal bound $N_\mathcal{T} = 8{,}190$ is not tight
here; it becomes tight in the exploitation regime ($T \gg N_\mathcal{T}$).
A tighter bound matching the empirical rate would require analysing UCT's traversal
concentration, which we leave to future work.
\end{proof}

\subsection{Proof of Theorem~\ref{thm:lower}}
\label{app:lower}

\begin{proof}
We bound \emph{cumulative} regret $R(T) = \sum_{t=1}^T [V(s_t^*) - V(\hat{s}_t^*)]$
by lower-bounding the per-step error probability at every step $t$, not merely at
the final step.

\paragraph{Hard instance family.}
Let $\mathcal{T}$ consist of root $s_0$ with two children $s_+$ and $s_-$.
Choose $\Delta = \sigma/(2\sqrt{T})$ and let $\theta \in \{+1,-1\}$ be chosen
uniformly at random, defining:
\begin{align*}
  \mathcal{I}^{(+)}: &\quad V(s_+) = \tfrac{1}{2}+\Delta,\; V(s_-) = \tfrac{1}{2}-\Delta,\\
  \mathcal{I}^{(-)}: &\quad V(s_+) = \tfrac{1}{2}-\Delta,\; V(s_-) = \tfrac{1}{2}+\Delta.
\end{align*}
Verifier observations satisfy $\hat{V}(s) = V(s)+\varepsilon$ with
$\varepsilon \sim \mathcal{N}(0,\sigma^2)$.
For $t \geq 2$ (after both nodes have been visited at least once),
$s_t^* = s_+$ under $\mathcal{I}^{(+)}$ and $s_t^* = s_-$ under $\mathcal{I}^{(-)}$.

\paragraph{Per-step regret via total variation.}
Denote by $\mathcal{H}_t$ the algorithm's complete history up to step $t$ and let
$n_t = \sum_{i=1}^t b_i$ be the total verifier calls through step $t$.
Each call contributes KL divergence $\mathrm{KL}(\mathcal{N}(\Delta,\sigma^2)\|\mathcal{N}(-\Delta,\sigma^2))
= 2\Delta^2/\sigma^2$ between the two instances.
By the chain rule and Pinsker's inequality:
\[
  \mathrm{TV}\!\left(\mathcal{H}_t^{(+)},\, \mathcal{H}_t^{(-)}\right)
  \;\leq\; \sqrt{\tfrac{1}{2}\,n_t \cdot \tfrac{2\Delta^2}{\sigma^2}}
  \;=\; \Delta\sqrt{\tfrac{n_t}{\sigma^2}}.
\]
Since $\hat{s}_t^*$ is a measurable function of $\mathcal{H}_t$, the
data-processing inequality and Le~Cam's two-point method give:
\[
  \Pr\!\left[\hat{s}_t^* \neq s_t^*\right]
  \;\geq\; \tfrac{1}{2}\!\left(1 - \mathrm{TV}\!\left(\mathcal{H}_t^{(+)},\mathcal{H}_t^{(-)}\right)\right)
  \;\geq\; \tfrac{1}{2}\!\left(1 - \Delta\sqrt{n_t/\sigma^2}\right).
\]
When $\hat{s}_t^* \neq s_t^*$, the per-step regret is $V(s_t^*) - V(\hat{s}_t^*) = 2\Delta$.
Hence:
\[
  \mathbb{E}[r_t] \;=\; 2\Delta\,\Pr\!\left[\hat{s}_t^* \neq s_t^*\right]
  \;\geq\; \Delta\!\left(1 - \Delta\sqrt{n_t/\sigma^2}\right).
\]

\paragraph{Summing over all steps.}
Let $B_T = n_T = \sum_{t=1}^T b_t$ be the total verifier calls.
Since $n_t \leq B_T$ for all $t$:
\begin{align}
  \mathbb{E}[R(T)]
  &= \sum_{t=2}^T \mathbb{E}[r_t]
   \;\geq\; (T-1)\Delta\!\left(1 - \Delta\sqrt{B_T/\sigma^2}\right). \label{eq:lb_sum}
\end{align}

\paragraph{Setting $\Delta$ and concluding.}
Substituting $\Delta = \sigma/(2\sqrt{T})$:
\[
  \Delta\sqrt{B_T/\sigma^2} = \frac{\sqrt{B_T}}{2\sqrt{T}} \;\leq\; \frac{1}{2},
\]
where the last inequality holds because $B_T \leq T$ (at most one call per step).
Inserting into \eqref{eq:lb_sum}:
\[
  \mathbb{E}[R(T)]
  \;\geq\; (T-1) \cdot \frac{\sigma}{2\sqrt{T}} \cdot \left(1 - \tfrac{1}{2}\right)
  \;=\; \frac{\sigma(T-1)}{4\sqrt{T}}
  \;\geq\; \frac{\sigma}{8}\sqrt{T}
  \qquad\text{for all } T \geq 2.
\]
This bound holds for \emph{every} policy $\pi$ on this instance family,
regardless of the number of verifier calls made.
Hence $\mathbb{E}[R(T)] \geq \frac{\sigma}{8}\sqrt{T}$.
\end{proof}

\section{Hyperparameter Sensitivity}
\label{app:hyperparams}

Table~\ref{tab:hyperparams} reports the sensitivity of DVA-MCTS to the threshold
exponent $\alpha$ with $\gamma = 1.0$ fixed, from the ablation experiment
(Section~\ref{sec:ablation}, 50 runs, $T = 400$, $L = 0.35$, $\sigma = 0.05$).

\begin{table}[h]
\centering
\caption{Ablation over threshold exponent $\alpha$ at $T{=}400$, $\gamma{=}1.0$
         (50 runs; all numbers from \texttt{results/ablation\_threshold.json}).}
\label{tab:hyperparams}
\begin{tabular}{cccc}
\toprule
$\alpha$ & Regret $R(400)$ & Verifier Calls & Call Fraction \\
\midrule
$0.25$ & $27.87$ & $371.9$ & $93.0\%$ \\
$0.50$ & $18.98$ & $382.2$ & $95.5\%$ \\
$0.75$ & $\mathbf{17.82}$ & $382.7$ & $95.7\%$ \\
$1.00$ & $19.99$ & $387.0$ & $96.7\%$ \\
\bottomrule
\end{tabular}
\end{table}

The empirically optimal exponent at $T{=}400$ is $\alpha{=}0.75$ (lowest regret).
The theoretically recommended $\alpha{=}0.5$ is within $6.5\%$ of optimal regret.
As discussed in Section~\ref{sec:ablation}, verification rates are uniformly high
at $T{=}400$ because the tree remains largely unexplored; the regime where $\alpha$
most strongly differentiates call patterns is large $B$ (${\geq}800$).

\section{Notation Summary}
\label{app:notation}

\begin{table}[h]
\centering
\caption{Summary of notation.}
\label{tab:notation}
\begin{tabular}{cl}
\toprule
Symbol & Description \\
\midrule
$\mathcal{T}$ & Search tree \\
$\mathcal{S}, \mathcal{A}$ & State and action spaces \\
$K, D$ & Branching factor and depth \\
$V(s)$ & True verifier score at state $s$ \\
$\hat{V}(s)$ & Noisy verifier observation at $s$ \\
$\hat{v}_t$ & Algorithm's estimate at step $t$ \\
$b_t \in \{0,1\}$ & Verifier call indicator at step $t$ \\
$\tau_t$ & Budget-sensitive threshold at step $t$ \\
$\gamma$ & Threshold constant \\
$\alpha$ & Threshold decay exponent \\
$L$ & Lipschitz constant of $V$ \\
$\sigma$ & Sub-Gaussian noise parameter \\
$T$ & Total search budget \\
$R(T)$ & Cumulative regret \\
$\mathcal{B}(\pi)$ & Total verifier calls under policy $\pi$ \\
$s^*_t$ & Oracle's best state through step $t$ \\
$\hat{s}^*_t$ & Algorithm's best estimated state through step $t$ \\
$s_{\text{anc}}$ & Deepest ancestor with stored verifier observation \\
$\mathcal{M}$ & Score memory (set of $(s, \hat{V}(s))$ pairs) \\
\bottomrule
\end{tabular}
\end{table}
